% !PRINTED 242
\weberpage{242}
%<<CONT
belongs. Since, together with the derivatives $f'(\xi+i\eta), f''(\xi+i\eta)\ldots$, the conjugate quantities $f'(\xi-i\eta), f''(\xi-i\eta)\ldots$ vanish at the same time, it follows that \gesperrt{conjugate imaginary roots have the same degree of multiplicity.}

From this it follows that the imaginary roots always occur in even number, and that an equation of odd degree must always have at least one real root.

Our next aim will be to determine the number of real roots, without solving the equation, directly from the values of \gesperrt{certain rational functions of the coefficients of the equation}.

A very important part is played here by the discriminant, whose significance for our question we must first consider.

In \S.\,46 we defined the discriminant as the product of the squares of all the root differences
\[
(x_1 - x_2)^2,
\]
further multiplied by $a_0^{2n-2}$, where $a_0$ is the coefficient of the highest power of the unknown in the given equation, and we showed there how the discriminant can be computed as a rational function of the coefficients. The factor $a_0^{2n-2}$ is always positive, and could here also be assumed equal to 1 without essential restriction of generality. As before, we shall use the sign $D$ for the discriminant.

\gesperrt{The discriminant vanishes when, and only when, two equal roots occur among the roots.}

If we assume that $f(x)$ has been freed of multiple factors by separating off the greatest common divisor of $f(x)$ and $f'(x)$, which is done by rational computation, then $D$ will not vanish.

If $x_1$ and $x_2$ are real, then $(x_1 - x_2)^2$ is positive. If $x_1$ is real and $x_2$ imaginary, then there is a root $x_2'$ conjugate to $x_2$, and the product
\[
(x_1 - x_2)(x_1 - x_2')
\]
is positive, being the product of two conjugate imaginary quantities, and hence so is its square.
